\documentclass[11pt]{article}
\usepackage{fontspec,amsmath,amssymb,amsthm,hyperref}
\usepackage[margin=1in]{geometry}
\setlength{\emergencystretch}{3em}
\title{The exact potential and average smoothness constants}
\author{LangevinLean mathematical exposition}
\date{}
\begin{document}\maketitle
\noindent Source module: \texttt{AverageLMC/Assumptions.lean}. Compile this proof document with XeLaTeX or LuaLaTeX. The accompanying HTML pairs the mathematical argument with the exact Lean statements and records the scope of the formalization.


This module represents the paper's state space by EuclideanSpace over Fin p. Every derivative and smoothness constant is computed from one real-valued potential. Coordinate smoothness uses the literal supremum of the Taylor quotients in (coordM); it is not a supplied error bound. The deterministic results assume differentiability and a Lipschitz gradient, without adding a Hessian or third derivative assumption.

\[M_j=\sup_{x,\ a\ne0}\frac{2(f(x+ae_j)-f(x)-a\,\partial_j f(x))}{a^2},\qquad M_{\rm av}=\frac1p\sum_jM_j.\]
\section{Euclidean coordinates and literal Taylor quotients}


\texttt{State p} is the real Euclidean space of dimension p and \texttt{coordinateVector j} is its jth unit vector. The coordinate Taylor quotient is twice the difference between f(x+a e\_j) and its first order approximation at x, divided by a squared. The defining set ranges over every x and every nonzero real a. Its supremum is exactly M\_j, and the average of these p constants is M\_av. A displacement a=1 at x=0 proves that every coordinate quotient set is nonempty. The unit norm and coordinate inner-product identities connect these definitions to Euclidean calculus.



\noindent Lean declarations: \texttt{State}, \texttt{coordinateVector}, \texttt{coordinateTaylorQuotient}, \texttt{coordinateTaylorQuotients}, \texttt{coordinateSmoothness}, \texttt{averageSmoothness}, \texttt{coordinateTaylorQuotients\_nonempty}, \texttt{norm\_coordinateVector}, \texttt{inner\_coordinateVector}.

\section{The actual global gradient seminorm}


\texttt{globalSmoothness} is the supremum of the gradient difference norm divided by the state difference norm over unequal states. For p positive, a pair consisting of zero and one coordinate vector witnesses nonemptiness. Any supplied gradient Lipschitz constant M bounds every quotient, so it bounds the supremum. The quotients are nonnegative, hence the supremum is nonnegative. Conversely, each quotient is at most this bounded supremum. Multiplying by the strictly positive state difference norm proves that the gradient is Lipschitz with its exact seminorm. Thus subsequent inequalities may use the literal constant in the paper instead of an arbitrary larger bound.



\noindent Lean declarations: \texttt{globalSmoothness}, \texttt{globalSmoothness\_le\_bound}, \texttt{globalSmoothness\_nonneg}, \texttt{gradient\_lipschitz\_globalSmoothness}.

\section{Bound and attain the coordinate quadratic model}


The shared differentiable descent lemma gives an upper quadratic model with coefficient M/2. At a coordinate displacement a e\_j, the inner product is a times the jth gradient coordinate and the squared displacement norm is a squared. For nonzero a, multiply the upper model by 2 and divide by a squared to show that every Taylor quotient is at most M. This proves boundedness of the quotient set and M\_j at most M. Each quotient is also at most its supremum M\_j. Multiplication by a squared and rearrangement recover the sharper coordinate upper model with coefficient M\_j/2. The zero displacement follows directly. No existence of a second derivative is needed.



\noindent Lean declarations: \texttt{coordinateTaylorQuotient\_le}, \texttt{coordinateTaylorQuotients\_bddAbove}, \texttt{coordinateSmoothness\_le\_global\_bound}, \texttt{coordinate\_upper\_bound}.

\section{Compare strong convexity, average smoothness and global smoothness}


The strong-convexity lower model at zero and e\_j gives the same first order term as the coordinate upper model, with quadratic coefficients m/2 and M\_j/2 respectively. Subtraction yields m at most M\_j. Summing the upper and lower bounds over the p coordinates and dividing by positive p proves m at most M\_av at most M. Apply the upper bounds to the actual nonnegative global gradient seminorm, using its proved Lipschitz property, to obtain M\_j at most M\_infinity and M\_av at most M\_infinity with exactly the paper's supremum constants.



\noindent Lean declarations: \texttt{strongConvex\_le\_coordinateSmoothness}, \texttt{averageSmoothness\_le\_global\_bound}, \texttt{coordinateSmoothness\_le\_globalSmoothness}, \texttt{averageSmoothness\_le\_globalSmoothness}, \texttt{strongConvex\_le\_averageSmoothness}.

\section{Actual Hessian, Laplacian and finite-sum notation}


\texttt{hessian} is the second Frechet derivative evaluated on pairs of coordinate vectors, and \texttt{laplacian} is mathlib's Laplacian. Expanding the Laplacian in the canonical orthonormal basis gives exactly the matrix trace of this Hessian. \texttt{LaplacianSmooth} records three continuous derivatives and the displayed bound on the squared norm of the gradient of the Laplacian, with the factor p from (lip:2). It does not introduce a Hessian Lipschitz requirement. Finally, \texttt{rmsComponentSmoothness} is the square root of the average of the squared exact component-gradient seminorms, as in (eq:Mtwo).



\noindent Lean declarations: \texttt{hessian}, \texttt{laplacian}, \texttt{LaplacianSmooth}, \texttt{rmsComponentSmoothness}, \texttt{trace\_hessian}.

\end{document}
